% ==============================================================================
% CHAPITRE 3 : IMPLÉMENTATION
% ==============================================================================
\chapter{Implémentation}

\begin{infobox}[Objectif du chapitre]
Ce chapitre décrit en détail l'environnement de développement, la structure
du projet, les algorithmes clés et les parties importantes du code source
de chaque composant du système.
\end{infobox}

% ──────────────────────────────────────────────────────────────────────────────
\section{Environnements de Développement}
% ──────────────────────────────────────────────────────────────────────────────

\subsection{Environnement matériel}

\begin{table}[H]
\centering
\caption{Configuration matérielle de développement}
\label{tab:hardware}
\renewcommand{\arraystretch}{1.4}
\begin{tabular}{lll}
\toprule
\rowcolor{bleuprincipal!15}
\textbf{Composant} & \textbf{Spécification} & \textbf{Impact} \\
\midrule
Processeur & Intel/AMD x86-64 & Exécution du modèle sur CPU \\
\rowcolor{gray!8}
Mémoire RAM & 8 Go minimum & Chargement du modèle ($\sim$430 Mo) \\
Stockage & 20 Go disponibles & Modèle + SDK Android + venv \\
\rowcolor{gray!8}
OS & Ubuntu 24.04 LTS & Environnement de production Linux \\
Réseau & Wi-Fi / Ethernet & Communication émulateur $\leftrightarrow$ API \\
\bottomrule
\end{tabular}
\end{table}

\subsection{Environnement logiciel}

\begin{table}[H]
\centering
\caption{Logiciels et versions utilisés}
\label{tab:software}
\renewcommand{\arraystretch}{1.4}
\begin{tabular}{lll}
\toprule
\rowcolor{bleuprincipal!15}
\textbf{Logiciel} & \textbf{Version} & \textbf{Rôle} \\
\midrule
Python & 3.12 & Langage backend et IA \\
\rowcolor{gray!8}
PyTorch & 2.x & Framework Deep Learning \\
Transformers (HuggingFace) & 4.x & Modèles pré-entraînés \\
\rowcolor{gray!8}
FastAPI & 0.110+ & Framework API REST \\
SQLAlchemy & 2.0+ & ORM Python-PostgreSQL \\
\rowcolor{gray!8}
PostgreSQL & 16 & Base de données relationnelle \\
Streamlit & 1.x & Dashboard web Python \\
\rowcolor{gray!8}
Flutter SDK & 3.24.5 & Développement application mobile \\
Dart & 3.x & Langage de l'application mobile \\
\rowcolor{gray!8}
Android Studio & 2024.2 & Émulateur Android, IDE Flutter \\
VS Code & 1.119 & Éditeur de code principal \\
\rowcolor{gray!8}
Git & 2.x & Contrôle de version \\
\bottomrule
\end{tabular}
\end{table}

% ──────────────────────────────────────────────────────────────────────────────
\section{Structure du Projet}
% ──────────────────────────────────────────────────────────────────────────────

\begin{lstlisting}[
  language=bash, caption={Structure complète du projet},
  label={lst:structure}, basicstyle=\ttfamily\footnotesize,
  frame=single, backgroundcolor=\color{griscode},
  numbers=none, xleftmargin=0.5cm
]
sentiment_etudiants/
|-- ai_model/                  # Module IA (CamemBERT)
|   |-- preprocessing/
|   |   |-- nettoyage_texte.py # Pipeline de nettoyage NLP
|   |   |-- tokenisation.py    # Tokenisation CamemBERT
|   |   `-- preparer_donnees.py# Pipeline preparation dataset
|   |-- modele_camembert.py    # Chargement/sauvegarde du modele
|   |-- entrainement.py        # Script d'entrainement (fine-tuning)
|   |-- evaluation.py          # Metriques: accuracy, F1, matrice
|   `-- saved_model/           # Modele fine-tune sauvegarde
|-- backend/                   # API FastAPI
|   |-- database.py            # Connexion PostgreSQL via SQLAlchemy
|   |-- models/modeles_db.py   # ORM: table commentaires
|   |-- routers/commentaires.py# Routes REST: predict, CRUD
|   |-- services/prediction_service.py # Singleton inference
|   `-- main.py                # Point d'entree FastAPI
|-- dashboard/                 # Tableau de bord Streamlit
|   |-- app.py                 # Page d'accueil
|   |-- components/api_client.py# Client HTTP vers l'API
|   `-- pages/                 # Pages Streamlit multi-pages
|       |-- 1_Analyse.py       # Analyse individuelle
|       |-- 2_Statistiques.py  # Graphiques Plotly
|       `-- 3_Historique.py    # Liste et filtres
|-- mobile_app/                # Application Flutter
|   `-- lib/
|       |-- main.dart          # Point d'entree + navigation
|       |-- models/commentaire.dart  # Modele de donnees Dart
|       |-- services/api_service.dart# Client HTTP Dart
|       |-- widgets/sentiment_card.dart # Widget reutilisable
|       `-- screens/
|           |-- analyse_screen.dart  # Ecran d'analyse
|           `-- historique_screen.dart# Ecran historique
`-- datasets/
    |-- raw/commentaires_etudiants.csv # 225 commentaires
    `-- processed/{train,validation,test}.csv
\end{lstlisting}

% ──────────────────────────────────────────────────────────────────────────────
\section{Module IA — Pipeline CamemBERT}
% ──────────────────────────────────────────────────────────────────────────────

\subsection{Prétraitement du texte}

Avant toute tokenisation, les commentaires bruts sont nettoyés par un
pipeline séquentiel défini dans \texttt{nettoyage\_texte.py} :

\begin{lstlisting}[style=stylePython, caption={Pipeline de nettoyage NLP}, label={lst:nettoyage}]
def nettoyer_commentaire(texte: str) -> str:
    """Pipeline complet de nettoyage d'un commentaire."""
    if not isinstance(texte, str) or len(texte.strip()) == 0:
        return ""
    # Etape 1 : Normalisation Unicode (NFC)
    texte = unicodedata.normalize("NFC", texte)
    texte = texte.replace(u"«", '"').replace(u"»", '"')
    # Etape 2 : Suppression des URLs, emails et caracteres speciaux
    texte = re.sub(r"http\S+|www\S+", "", texte)
    texte = re.sub(r"\S+@\S+", "", texte)
    texte = re.sub(r"[^\w\sA-Za-zÀ-ž'',;:!?.\-]", " ", texte)
    # Etape 3 : Mise en minuscules
    texte = texte.lower()
    # Etape 4 : Normalisation des espaces
    texte = re.sub(r"\s+", " ", texte).strip()
    return texte
\end{lstlisting}

Ce pipeline traite les cas suivants :
\begin{itemize}
  \item \textbf{Normalisation Unicode :} harmonise les caractères composés
        (é peut être représenté de deux façons en Unicode).
  \item \textbf{Guillemets typographiques :} « » sont remplacés par des
        guillemets droits ASCII.
  \item \textbf{URLs et emails :} supprimés car non pertinents pour l'analyse
        de sentiment.
  \item \textbf{Minuscules :} CamemBERT est sensible à la casse, la mise
        en minuscules améliore la cohérence.
\end{itemize}

\subsection{Jeu de données}

Le jeu de données a été constitué manuellement avec \textbf{225 commentaires
étudiants} en français, répartis équitablement en 3 classes :

\begin{table}[H]
\centering
\caption{Répartition du jeu de données}
\label{tab:dataset}
\renewcommand{\arraystretch}{1.4}
\begin{tabular}{lccc}
\toprule
\rowcolor{bleuprincipal!15}
\textbf{Classe} & \textbf{Étiquette} & \textbf{Exemples} & \textbf{Proportion} \\
\midrule
Négatif & 0 & 75 & 33,33\% \\
\rowcolor{gray!8}
Neutre & 1 & 75 & 33,33\% \\
Positif & 2 & 75 & 33,33\% \\
\midrule
\rowcolor{bleuprincipal!8}
\textbf{Total} & -- & \textbf{225} & \textbf{100\%} \\
\bottomrule
\end{tabular}
\end{table}

Le découpage train/validation/test respecte une stratification :

\begin{table}[H]
\centering
\caption{Découpage du dataset (stratifié)}
\label{tab:split}
\renewcommand{\arraystretch}{1.4}
\begin{tabular}{lcccc}
\toprule
\rowcolor{bleuprincipal!15}
\textbf{Ensemble} & \textbf{Taille} & \textbf{Proportion} &
\textbf{Usage} & \textbf{Batches} \\
\midrule
Entraînement & 157 & $\sim$70\% & Mise à jour des poids & 10 \\
\rowcolor{gray!8}
Validation & 34 & $\sim$15\% & Sélection du meilleur modèle & 3 \\
Test & 34 & $\sim$15\% & Évaluation finale non biaisée & 3 \\
\bottomrule
\end{tabular}
\end{table}

\subsection{Architecture du modèle fine-tuné}

Le modèle CamemBERT fine-tuné est composé de :

\begin{enumerate}
  \item \textbf{Couche d'entrée :} tokenisation SentencePiece (vocabulaire
        32 000 tokens, séquences de max 128 tokens).
  \item \textbf{12 blocs Transformer :} chaque bloc contient un mécanisme
        d'attention multi-têtes (12 têtes, dimension 768) et un réseau
        Feed-Forward (dimension 3072).
  \item \textbf{Couche de classification :} projection linéaire
        $768 \rightarrow 3$ suivie d'un Softmax.
\end{enumerate}

Le nombre total de paramètres est de \textbf{110,6 millions}, dont les
poids de la couche de classification ($768 \times 3 + 3 = 2307$ paramètres
ajoutés).

\subsection{Boucle d'entraînement}

\begin{lstlisting}[style=stylePython, caption={Boucle d'entraînement fine-tuning CamemBERT}, label={lst:training}]
def entrainer_une_epoch(modele, loader_train, optimiseur, scheduler,
                        device, num_epoch):
    modele.train()
    perte_totale = 0.0
    nb_correct, nb_total = 0, 0

    for i, batch in enumerate(loader_train):
        input_ids      = batch["input_ids"].to(device)
        attention_mask = batch["attention_mask"].to(device)
        labels         = batch["labels"].to(device)

        # 1. Remise a zero des gradients
        optimiseur.zero_grad()

        # 2. Passage avant : calcul de la perte cross-entropique
        sorties = modele(input_ids=input_ids,
                         attention_mask=attention_mask,
                         labels=labels)
        perte = sorties.loss

        # 3. Retropropagation : calcul des gradients
        perte.backward()

        # 4. Ecrêtage des gradients (previent l'explosion)
        torch.nn.utils.clip_grad_norm_(modele.parameters(), max_norm=1.0)

        # 5. Mise a jour des poids
        optimiseur.step()
        scheduler.step()  # Ajuste le learning rate

        perte_totale += perte.item()
        predictions = torch.argmax(sorties.logits, dim=1)
        nb_correct  += (predictions == labels).sum().item()
        nb_total    += labels.size(0)

    return perte_totale / len(loader_train), nb_correct / nb_total * 100
\end{lstlisting}

\textbf{Points techniques importants :}
\begin{itemize}
  \item \texttt{optimiseur.zero\_grad()} : réinitialise les gradients
        accumulés du batch précédent. Sans cela, les gradients s'accumulent
        à chaque batch, causant des mises à jour erronées.
  \item \texttt{clip\_grad\_norm\_(max\_norm=1.0)} : technique standard pour
        les Transformers évitant l'explosion des gradients (gradient clipping).
  \item \texttt{scheduler.step()} : ajuste le taux d'apprentissage selon
        le scheduler linéaire avec warmup défini.
\end{itemize}

\subsection{Hyperparamètres d'entraînement}

\begin{table}[H]
\centering
\caption{Hyperparamètres du fine-tuning}
\label{tab:hyperparams}
\renewcommand{\arraystretch}{1.4}
\begin{tabular}{lll}
\toprule
\rowcolor{bleuprincipal!15}
\textbf{Hyperparamètre} & \textbf{Valeur} & \textbf{Justification} \\
\midrule
Taux d'apprentissage & $2 \times 10^{-5}$ & Standard fine-tuning BERT, évite le catastrophic forgetting \\
\rowcolor{gray!8}
Taille de lot (batch) & 16 & Compromis mémoire/variance du gradient \\
Époques & 3 & Suffisant pour un petit dataset, évite le surapprentissage \\
\rowcolor{gray!8}
Max tokens & 128 & Couvre 95\%+ des commentaires, économise la mémoire \\
Optimiseur & AdamW & Standard pour les Transformers, weight decay régularise \\
\rowcolor{gray!8}
Warmup ratio & 10\% & Stabilise le début de l'entraînement \\
Gradient clipping & 1.0 & Prévient l'explosion des gradients \\
\rowcolor{gray!8}
Weight decay & 0.01 & Régularisation L2 pour réduire le surapprentissage \\
\bottomrule
\end{tabular}
\end{table}

% ──────────────────────────────────────────────────────────────────────────────
\section{Module Backend — FastAPI}
% ──────────────────────────────────────────────────────────────────────────────

\subsection{Démarrage du serveur et lifespan}

\begin{lstlisting}[style=stylePython, caption={Initialisation FastAPI avec lifespan}, label={lst:fastapi_main}]
@asynccontextmanager
async def lifespan(app: FastAPI):
    """Initialisation au demarrage, nettoyage a l'arret."""
    print("[DEMARRAGE] Initialisation de l'API...")
    creer_tables()           # Creation des tables PostgreSQL
    service_prediction.charger()  # Chargement du modele CamemBERT en RAM
    print("[DEMARRAGE] API prete.")
    yield                    # Le serveur tourne ici
    print("[ARRET] Fermeture de l'API.")

app = FastAPI(
    title="API Analyse de Sentiments Etudiants",
    version="1.0.0",
    lifespan=lifespan
)

app.add_middleware(CORSMiddleware,
    allow_origins=["*"],  # Flutter + Streamlit peuvent appeler l'API
    allow_methods=["*"],
    allow_headers=["*"]
)
\end{lstlisting}

Le pattern \texttt{lifespan} remplace les anciens \texttt{startup\_event}
et \texttt{shutdown\_event} de FastAPI. Il garantit que le modèle CamemBERT
est chargé une seule fois en mémoire au démarrage du serveur, évitant un
rechargement coûteux à chaque requête.

\subsection{Service de prédiction — Pattern Singleton}

\begin{lstlisting}[style=stylePython, caption={Implémentation du pattern Singleton}, label={lst:singleton}]
class ServicePrediction:
    """
    Singleton : une seule instance partagee dans toute l'application.
    Garantit que le modele CamemBERT (430 Mo) n'est charge qu'une fois.
    """
    _instance = None

    def __new__(cls):
        if cls._instance is None:
            cls._instance = super().__new__(cls)
            cls._instance._initialise = False
        return cls._instance

    def predict(self, texte: str) -> dict:
        texte_nettoye = nettoyer_commentaire(texte)

        encodage = self.tokeniseur(
            texte_nettoye,
            max_length=128,
            padding="max_length",
            truncation=True,
            return_tensors="pt"
        )

        with torch.no_grad():  # Desactive le calcul des gradients
            sorties = self.modele(
                input_ids=encodage["input_ids"].to(self.device),
                attention_mask=encodage["attention_mask"].to(self.device)
            )

        # Conversion logits -> probabilites
        probs = F.softmax(sorties.logits, dim=1).squeeze().tolist()
        classe = int(torch.argmax(sorties.logits, dim=1).item())

        return {
            "sentiment_predit": classe,
            "label_sentiment" : LABELS[classe],
            "score_negatif"   : round(probs[0], 4),
            "score_neutre"    : round(probs[1], 4),
            "score_positif"   : round(probs[2], 4),
        }
\end{lstlisting}

\textbf{Pourquoi \texttt{torch.no\_grad()} ?} Lors de l'inférence (prédiction),
nous n'avons pas besoin de calculer les gradients (ça, c'est pour
l'entraînement). Cette directive réduit la consommation mémoire de 50\%
et accélère les calculs.

\subsection{Validation des données avec Pydantic}

\begin{lstlisting}[style=stylePython, caption={Schémas de validation Pydantic}, label={lst:pydantic}]
class RequetePrediction(BaseModel):
    """Corps de la requete POST /api/predict"""
    texte   : str           = Field(..., min_length=5, max_length=2000,
                              description="Commentaire a analyser")
    matiere : Optional[str] = Field(None, max_length=100)

class ReponseCommentaire(BaseModel):
    """Schema de reponse standardise"""
    id               : int
    texte_original   : str
    matiere          : Optional[str]
    sentiment_predit : int
    label_sentiment  : str
    scores           : dict
    date_creation    : Optional[str]

    class Config:
        from_attributes = True  # Lecture depuis un objet SQLAlchemy
\end{lstlisting}

Pydantic garantit que :
\begin{itemize}
  \item Le texte fait au moins 5 caractères (sinon HTTP 422).
  \item Les types sont corrects (int, str, etc.).
  \item Les champs optionnels sont traités correctement.
  \item La sérialisation JSON est automatique.
\end{itemize}

\subsection{Routes API — Récapitulatif}

\begin{table}[H]
\centering
\caption{Endpoints de l'API REST}
\label{tab:endpoints}
\renewcommand{\arraystretch}{1.5}
\begin{tabular}{llp{3.5cm}p{3.5cm}}
\toprule
\rowcolor{bleuprincipal!15}
\textbf{Méthode} & \textbf{Route} & \textbf{Paramètres} & \textbf{Description} \\
\midrule
\rowcolor{vertprincipal!10}
POST & \texttt{/api/predict} & \texttt{\{texte, matiere?\}} & Analyse + sauvegarde \\
GET & \texttt{/api/commentaires} & \texttt{?limite, sentiment, matiere} & Liste paginée filtrée \\
\rowcolor{gray!8}
GET & \texttt{/api/commentaires/\{id\}} & \texttt{id} (chemin) & Détail d'un commentaire \\
GET & \texttt{/api/stats} & -- & Statistiques agrégées \\
\rowcolor{rougeprincipal!10}
DELETE & \texttt{/api/commentaires/\{id\}} & \texttt{id} (chemin) & Suppression \\
GET & \texttt{/health} & -- & Statut du serveur \\
\bottomrule
\end{tabular}
\end{table}

% ──────────────────────────────────────────────────────────────────────────────
\section{Module Application Mobile — Flutter}
% ──────────────────────────────────────────────────────────────────────────────

\subsection{Service API — Communication HTTP}

\begin{lstlisting}[style=styleDart, caption={Service HTTP Flutter}, label={lst:flutter_api}]
class ApiService {
  static const String _baseUrl = 'http://10.0.2.2:8000/api';
  static const Duration _timeout = Duration(seconds: 30);

  static Future<Commentaire?> analyserCommentaire({
    required String texte,
    String? matiere,
  }) async {
    try {
      final body = <String, dynamic>{'texte': texte};
      if (matiere != null) body['matiere'] = matiere;

      final reponse = await http.post(
        Uri.parse('$_baseUrl/predict'),
        headers: {'Content-Type': 'application/json; charset=utf-8'},
        body: jsonEncode(body),
      ).timeout(_timeout);

      if (reponse.statusCode == 200) {
        final json = jsonDecode(utf8.decode(reponse.bodyBytes));
        return Commentaire.fromJson(json);
      }
      return null;
    } catch (_) {
      return null;  // Gestion silencieuse des erreurs reseau
    }
  }
}
\end{lstlisting}

\textbf{Pourquoi \texttt{10.0.2.2} ?} L'émulateur Android est une machine
virtuelle qui possède sa propre interface réseau. Pour accéder au
\texttt{localhost} de l'ordinateur hôte depuis l'émulateur, Android
redirige l'adresse \texttt{10.0.2.2} vers \texttt{127.0.0.1} de la machine
hôte. Sur un vrai téléphone, il faudrait utiliser l'IP locale du PC.

\subsection{Navigation par onglets avec IndexedStack}

\begin{lstlisting}[style=styleDart, caption={Navigation IndexedStack}, label={lst:navigation}]
class _AccueilPageState extends State<AccueilPage> {
  int _ongletActif = 0;

  final List<Widget> _ecrans = const [
    AnalyseScreen(),     // Index 0
    HistoriqueScreen(),  // Index 1
  ];

  @override
  Widget build(BuildContext context) {
    return Scaffold(
      body: IndexedStack(
        index: _ongletActif,
        children: _ecrans,
      ),
      bottomNavigationBar: NavigationBar(
        selectedIndex: _ongletActif,
        onDestinationSelected: (i) => setState(() => _ongletActif = i),
        // ...
      ),
    );
  }
}
\end{lstlisting}

\texttt{IndexedStack} maintient les deux écrans en vie simultanément en
mémoire. L'écran visible est déterminé par l'index. Ce choix architectural
préserve l'état de chaque écran lors de la navigation (la liste de l'historique
n'est pas rechargée à chaque changement d'onglet).

% ──────────────────────────────────────────────────────────────────────────────
\section{Gestion des Erreurs et Performances}
% ──────────────────────────────────────────────────────────────────────────────

\subsection{Gestion des erreurs}

\begin{table}[H]
\centering
\caption{Stratégie de gestion des erreurs}
\label{tab:errors}
\renewcommand{\arraystretch}{1.4}
\begin{tabular}{p{3cm} p{3.5cm} p{6cm}}
\toprule
\rowcolor{bleuprincipal!15}
\textbf{Couche} & \textbf{Erreur possible} & \textbf{Gestion} \\
\midrule
API FastAPI & Texte trop court & HTTP 422 (Unprocessable Entity) + message \\
\rowcolor{gray!8}
API FastAPI & Erreur modèle & HTTP 500 + message d'erreur technique \\
API FastAPI & ID inexistant & HTTP 404 (Not Found) \\
\rowcolor{gray!8}
Flutter & Timeout réseau & \texttt{null} retourné, message affiché \\
Flutter & Serveur injoignable & Message ``Impossible de joindre l'API'' \\
\rowcolor{gray!8}
Modèle IA & Texte vide après nettoyage & ValueError explicite levée \\
\bottomrule
\end{tabular}
\end{table}

\subsection{Optimisations de performance}

\begin{itemize}
  \item \textbf{Singleton ServicePrediction :} le modèle CamemBERT (430 Mo)
        est chargé une seule fois en mémoire au démarrage. Toutes les requêtes
        partagent la même instance, évitant un rechargement à chaque appel.

  \item \textbf{torch.no\_grad() :} désactive le calcul du graphe de
        gradient lors de l'inférence, réduisant la mémoire consommée de
        50\% et améliorant la vitesse de 20\%.

  \item \textbf{Modèle en mode eval() :} \texttt{model.eval()} désactive
        le Dropout et le BatchNorm (actifs seulement à l'entraînement),
        garantissant des prédictions déterministes et plus rapides.

  \item \textbf{Pagination API :} les listes de commentaires sont paginées
        (paramètre \texttt{limite}) pour éviter de charger de grandes quantités
        de données inutilement.

  \item \textbf{IndexedStack Flutter :} évite la reconstruction des widgets
        à chaque navigation, préservant l'état et améliorant la fluidité.
\end{itemize}

\section*{Conclusion du Chapitre}

L'implémentation du système repose sur une architecture propre, modulaire
et respectueuse des bonnes pratiques. Chaque composant a été développé de
manière indépendante, communiquant via des interfaces bien définies.
Le chapitre suivant évalue les performances du système et présente les
résultats obtenus.

\cleardoublepage
